\documentclass[conference]{IEEEtran}

% Codificación y configuración general
\usepackage[utf8]{inputenc}
\usepackage[T1]{fontenc}
\usepackage[spanish,es-nodecimaldot]{babel}

% Paquetes habituales para artículos IEEE
\usepackage{cite}
\usepackage{amsmath,amssymb,amsfonts}
\usepackage{graphicx}
\usepackage{textcomp}
\usepackage{xcolor}
\usepackage{url}
\usepackage{booktabs}
\usepackage{listings}
\usepackage{algorithm}
\usepackage{algorithmic}

% TikZ se conserva para que pueda habilitarse cuando sea necesario.
\usepackage{tikz}
\usetikzlibrary{shapes.geometric,positioning}

\def\UrlBreaks{\do\/\do-\do\_}

\lstset{
    basicstyle=\ttfamily\small,
    showstringspaces=false,
    breaklines=true,
    frame=single,
    numbers=left,
    numberstyle=\tiny\color{gray},
    inputencoding=utf8,
    extendedchars=true,
    literate={á}{{\'a}}1 {é}{{\'e}}1 {í}{{\'i}}1 {ó}{{\'o}}1
    {ú}{{\'u}}1 {ñ}{{\~n}}1 {ü}{{\"u}}1 {Á}{{\'A}}1 {É}{{\'E}}1
    {Í}{{\'I}}1 {Ó}{{\'O}}1 {Ú}{{\'U}}1 {Ñ}{{\~N}}1 {Ü}{{\"U}}1
}

% Descomenta y ajusta esta ruta si el documento utiliza imágenes.
% \graphicspath{{../img/}}

\title{Título del trabajo}

\author{
\IEEEauthorblockN{Leonardo Daniel Aviña Neri}
\IEEEauthorblockA{
Licenciatura en Ingeniería de Datos e Inteligencia Artificial\\
Universidad de Guanajuato\\
\texttt{ld.avinaneri@ugto.mx}
}
\and
\IEEEauthorblockN{Nombre Apellido 2}
\IEEEauthorblockA{
Programa académico\\
Universidad o institución\\
\texttt{correo2@ejemplo.com}
}
\and
\IEEEauthorblockN{Nombre Apellido 3}
\IEEEauthorblockA{
Programa académico\\
Universidad o institución\\
\texttt{correo3@ejemplo.com}
}
\and
\IEEEauthorblockN{Nombre Apellido 4}
\IEEEauthorblockA{
Programa académico\\
Universidad o institución\\
\texttt{correo4@ejemplo.com}
}
\and
\IEEEauthorblockN{Nombre Apellido 5}
\IEEEauthorblockA{
Programa académico\\
Universidad o institución\\
\texttt{correo5@ejemplo.com}
}
}

\begin{document}
\maketitle

\begin{abstract}
Lorem ipsum dolor sit amet, consectetur adipiscing elit. Integer tincidunt,
nisl at facilisis consequat, justo sem porta neque, vitae tincidunt massa
libero non ipsum. Este resumen debe presentar brevemente el problema, la
metodología, los resultados principales y la conclusión del trabajo.
\end{abstract}

\begin{IEEEkeywords}
palabra clave 1, palabra clave 2, metodología, resultados, análisis
\end{IEEEkeywords}

\section{Introducción}
Lorem ipsum dolor sit amet, consectetur adipiscing elit. Quisque consequat,
urna sit amet fermentum ullamcorper, justo mauris posuere risus, vitae
convallis lectus erat vel arcu. En esta sección se presenta el contexto del
problema, su importancia y la motivación del trabajo.

El objetivo general del estudio es describir una propuesta reproducible para
analizar el problema seleccionado. Los objetivos específicos pueden incluir
la identificación de variables, el diseño de una solución y la evaluación de
los resultados obtenidos.

\section{Marco teórico}
Lorem ipsum dolor sit amet, consectetur adipiscing elit. Suspendisse potenti.
Praesent eget lectus sed sem commodo consequat. Las definiciones, conceptos
fundamentales y trabajos relacionados deben colocarse en esta sección.

Una expresión matemática puede escribirse de la siguiente manera:
\begin{equation}
    f(x) = \sum_{i=1}^{n} w_i x_i,
\end{equation}
donde $w_i$ representa un peso y $x_i$ representa una variable de entrada.

\section{Metodología}
Lorem ipsum dolor sit amet, consectetur adipiscing elit. Donec vitae
sollicitudin mauris. Nam euismod, massa quis consectetur tincidunt, sem
libero tincidunt lectus, at commodo justo neque quis erat.

\subsection{Materiales y herramientas}
Indique aquí los datos, herramientas, lenguajes de programación, bibliotecas
y recursos empleados en el desarrollo.

\subsection{Procedimiento}
Describa aquí las etapas necesarias para reproducir el trabajo. Puede
organizar el procedimiento mediante una lista, un algoritmo o subsecciones.

\begin{algorithm}[htbp]
\caption{Procedimiento de ejemplo}
\label{alg:procedimiento}
\begin{algorithmic}[1]
\STATE Recibir los datos de entrada.
\STATE Procesar la información de acuerdo con el método propuesto.
\STATE Evaluar y registrar los resultados.
\RETURN Resultados del análisis.
\end{algorithmic}
\end{algorithm}

\section{Resultados}
Lorem ipsum dolor sit amet, consectetur adipiscing elit. Curabitur ac
sollicitudin justo. Los resultados deben presentarse de forma clara,
acompañados de tablas, figuras o métricas cuando sea necesario.

\begin{table}[htbp]
\centering
\caption{Tabla de resultados de ejemplo.}
\label{tab:resultados}
\begin{tabular}{lcc}
\toprule
Categoría & Métrica 1 & Métrica 2 \\
\midrule
Caso A & 0.00 & 0.00 \\
Caso B & 0.00 & 0.00 \\
\bottomrule
\end{tabular}
\end{table}

% Gráfico de ejemplo con TikZ. Descomenta este bloque cuando lo necesites.
% \begin{figure}[htbp]
% \centering
% \begin{tikzpicture}[
%     block/.style={
%         rectangle,
%         rounded corners=3pt,
%         draw=black,
%         align=center,
%         minimum width=0.75\columnwidth,
%         minimum height=0.7cm
%     },
%     arrow/.style={-stealth, draw=black},
%     node distance=0.4cm
% ]
% \node[block] (entrada) {Datos de entrada};
% \node[block, below=of entrada] (proceso) {Procesamiento};
% \node[block, below=of proceso] (salida) {Resultados};
% \draw[arrow] (entrada) -- (proceso);
% \draw[arrow] (proceso) -- (salida);
% \end{tikzpicture}
% \caption{Diagrama de flujo de ejemplo.}
% \label{fig:diagrama}
% \end{figure}

\section{Discusión}
Lorem ipsum dolor sit amet, consectetur adipiscing elit. Aenean sed
ullamcorper mi. Compare aquí los resultados con los objetivos planteados,
explique sus implicaciones y señale las limitaciones del estudio.

\section{Conclusiones}
Lorem ipsum dolor sit amet, consectetur adipiscing elit. A partir del
análisis realizado, redacte las conclusiones principales, las aportaciones
del trabajo y las posibles líneas de investigación futura.

\section*{Agradecimientos}
Opcionalmente, incluya aquí agradecimientos a personas, instituciones o
proyectos que hayan contribuido al desarrollo del trabajo.

\begin{thebibliography}{00}
\bibitem{referencia1}
Autor, ``Título de la referencia,'' \textit{Nombre de la revista o
conferencia}, vol. 1, no. 1, pp. 1--10, 2026.

\bibitem{referencia2}
Autor, \textit{Título del libro}, Editorial, 2026.
\end{thebibliography}

\end{document}
